%% Lecture 16 -- Magnetic dipoles in external fields and the beginning of magnetic materials
\section{Lecture 16 -- Magnetic Dipoles in External Fields and Magnetic Materials}\label{sec:lec16}

\begin{center}
\begin{tabular}{ll}
  \textbf{Date:}\quad 31/05/2026 & \\
\end{tabular}
\end{center}
\hrule\vspace{1.5em}

\subsection*{Overview}

Lecture 15 introduced the magnetic dipole moment
$\bmu$ as the leading far-field data of a localized current distribution.  This
lecture explains why that concept is the right microscopic language for magnetic
materials.  A piece of matter can be viewed, at sufficiently long wavelengths, as
many tiny current loops: orbital electron currents and spin magnetic moments.  The
external magnetic field does not see the detailed atomic motion directly; it sees
the dipole moments those motions produce.

The main results are:
\begin{itemize}
  \item the magnetic field of a point magnetic dipole is
        $\bB_{\rm dip}(\br)=\bigl[3\hat{\br}(\bmu\cdot\hat{\br})-\bmu\bigr]/r^3$
        away from the dipole;
  \item the interaction term for a small current loop in a slowly varying external
        field is $L_{\rm int}=\bmu\cdot\bB_{\rm ext}$;
  \item therefore the potential energy is $U=-\bmu\cdot\bB_{\rm ext}$;
  \item the force on a dipole is $\bvec{F}=\nabla(\bmu\cdot\bB)$, so a uniform
        magnetic field exerts no net force on an ideal dipole;
  \item the torque is $\boldsymbol{\tau}=\bmu\times\bB$, so a uniform field can
        rotate the dipole even when it cannot translate it;
  \item for orbital motion, combining $\bmu=(e/2m_ec)\bvec{L}$ with
        $d\bvec{L}/dt=\boldsymbol{\tau}$ gives Larmor precession;
  \item magnetic materials are classified by how their microscopic dipoles respond:
        paramagnets align with the applied field, diamagnets oppose it,
        ferromagnets retain magnetization, and superconductors expel magnetic field.
\end{itemize}

%------------------------------------------------------------
\subsection{Magnetic Field of a Dipole}\label{subsec:lec16-dipole-field}

\subsubsection*{Physical Motivation}
Last lecture gave the far-zone vector potential of a localized current distribution,
\begin{equation}\label{eq:lec16-Adip-start}
  \Avec_{\rm dip}(\br)=\frac{\bmu\times\br}{r^3}.
\end{equation}
For matter, this is the microscopic building block: an atom or molecule is replaced
by a pointlike magnetic dipole when we observe it on length scales much larger than
atomic size.  To understand how many such dipoles produce magnetization, we first need
the magnetic field of one dipole.

\subsubsection*{Derivation}
Use the vector identity
\begin{equation}\label{eq:lec16-curl-cross-identity}
  \nabla\times(\bvec{a}\times\bvec{c})
  =
  \bvec{a}(\nabla\cdot\bvec{c})
  -\bvec{c}(\nabla\cdot\bvec{a})
  +(\bvec{c}\cdot\nabla)\bvec{a}
  -(\bvec{a}\cdot\nabla)\bvec{c}.
\end{equation}
Here $\bvec{a}=\bmu$ is constant and $\bvec{c}=\br/r^3$, so
$\nabla\cdot\bmu=0$ and $(\br/r^3\cdot\nabla)\bmu=0$.  Thus
\begin{equation}\label{eq:lec16-Bdip-identity}
  \bB_{\rm dip}
  =\nabla\times\left(\frac{\bmu\times\br}{r^3}\right)
  =
  \bmu\,\nabla\cdot\left(\frac{\br}{r^3}\right)
  -(\bmu\cdot\nabla)\left(\frac{\br}{r^3}\right).
\end{equation}
Away from the origin,
\begin{equation}\label{eq:lec16-div-r-over-r3}
  \nabla\cdot\left(\frac{\br}{r^3}\right)=0
  \qquad (r\neq 0),
\end{equation}
while componentwise
\begin{align}
  (\bmu\cdot\nabla)\left(\frac{r_i}{r^3}\right)
  &= \mu_j\partial_j(r_i r^{-3}) \notag\\
  &= \mu_j\left(\delta_{ij}r^{-3}+r_i\,\partial_j r^{-3}\right) \notag\\
  &= \frac{\mu_i}{r^3}-3\,\frac{r_i(\bmu\cdot\br)}{r^5}. \label{eq:lec16-directional-derivative}
\end{align}
Substituting this into \eqref{eq:lec16-Bdip-identity} gives
\begin{equation}\label{eq:lec16-Bdip}
  \boxed{
    \bB_{\rm dip}(\br)
    =
    \frac{3\br(\bmu\cdot\br)}{r^5}
    -\frac{\bmu}{r^3}
    =
    \frac{3\hat{\br}(\bmu\cdot\hat{\br})-\bmu}{r^3},
    \qquad r\neq 0.
  }
\end{equation}
\textit{Significance:} The dipole field decays as $1/r^3$, exactly like the electric
dipole field.  The formula is meaningful only outside the microscopic source; if one
approaches the current loop itself, the point-dipole approximation must be replaced by
the actual current distribution.

\begin{remark}
  Distributionally, there is also a contact term at the dipole position.  In this
  lecture we use the far-field expression \eqref{eq:lec16-Bdip}, because magnetic
  materials are modeled by replacing each small current loop by its field outside
  atomic scales.
\end{remark}

\begin{figure}[h]
\centering
\begin{tikzpicture}[>=Stealth, thick, scale=1.05]
  \draw[->, gray] (-3,0) -- (3,0) node[right] {$x$};
  \draw[->, gray] (0,-2.5) -- (0,2.6) node[above] {$z$};
  \draw[->, red!70!black, line width=1.1pt] (0,-0.65) -- (0,0.9)
    node[right] {$\bmu$};
  \fill[black] (0,0) circle (2pt);

  \foreach \x in {-2.1,-1.4,-0.8,0.8,1.4,2.1} {
    \draw[blue!60!black, ->]
      (\x,-1.7) .. controls ({0.45*\x},-0.8) and ({0.45*\x},0.8) .. (\x,1.7);
  }
  \draw[blue!60!black, ->]
    (-0.35,2.05) .. controls (-1.0,1.75) and (-1.1,1.1) .. (-0.65,0.65);
  \draw[blue!60!black, ->]
    (0.35,2.05) .. controls (1.0,1.75) and (1.1,1.1) .. (0.65,0.65);
  \draw[blue!60!black, ->]
    (-0.35,-2.05) .. controls (-1.0,-1.75) and (-1.1,-1.1) .. (-0.65,-0.65);
  \draw[blue!60!black, ->]
    (0.35,-2.05) .. controls (1.0,-1.75) and (1.1,-1.1) .. (0.65,-0.65);
  \node[blue!60!black] at (2.25,2.1) {$\bB_{\rm dip}$};
\end{tikzpicture}
\caption{Schematic field pattern of a magnetic dipole.  Far from the microscopic
current loop, the entire source is represented by the single vector $\bmu$.}
\label{fig:lec16-dipole-field}
\end{figure}

%------------------------------------------------------------
\subsection{Interaction of a Small Current Loop with an External Field}\label{subsec:lec16-interaction}

\subsubsection*{Physical Motivation}
To study a material in an applied magnetic field, we need the energy of one microscopic
current loop in that field.  The applied field is assumed to vary slowly across the
loop, because the loop is atomic in size while laboratory fields vary on macroscopic
scales.  This is the same logic as a multipole expansion: keep the leading term first.

The covariant interaction between matter current and electromagnetic field can be
written as
\begin{equation}\label{eq:lec16-Sint-covariant}
  S_{\rm int}
  =
  -\frac{1}{c^2}\int J^\mu A_\mu\,d^4x.
\end{equation}
In the magnetostatic neutral sector, $J^0=0$ and $A^0=\phi=0$.  With the mostly-minus
metric convention used in this course, this reduces to
\begin{equation}\label{eq:lec16-Sint-current}
  S_{\rm int}
  =
  \frac{1}{c}\int dt\int \Jvec(\br)\cdot\Avec(\br)\,d^3r.
\end{equation}
For a thin current loop carrying current $I$,
\begin{equation}\label{eq:lec16-J-to-loop}
  \Jvec(\br)\,d^3r \longrightarrow I\,d\bvec{\ell},
\end{equation}
so
\begin{align}
  S_{\rm int}
  &= \frac{I}{c}\int dt\oint_C \Avec\cdot d\bvec{\ell} \notag\\
  &= \frac{I}{c}\int dt\iint_\Sigma
     (\nabla\times\Avec)\cdot d\bvec{S} \notag\\
  &= \frac{I}{c}\int dt\iint_\Sigma \bB\cdot d\bvec{S}. \label{eq:lec16-Sint-stokes}
\end{align}
If $\bB$ is nearly constant over the loop,
\begin{equation}\label{eq:lec16-uniform-field-loop}
  \iint_\Sigma \bB\cdot d\bvec{S}
  \simeq
  \bB\cdot\boldsymbol{\Sigma},
\end{equation}
where $\boldsymbol{\Sigma}$ is the oriented area vector.  Since
$\bmu=(I/c)\boldsymbol{\Sigma}$, we obtain
\begin{equation}\label{eq:lec16-Lint}
  \boxed{L_{\rm int}=\bmu\cdot\bB_{\rm ext}.}
\end{equation}
\textit{Significance:} The microscopic loop couples to the external field only through
its dipole moment at leading order.  This is why $\bmu$ is the natural variable for
magnetic matter.

Since a mechanical Lagrangian has the structure $L=T-U$, the potential energy is
\begin{equation}\label{eq:lec16-U-dipole}
  \boxed{U=-\bmu\cdot\bB_{\rm ext}.}
\end{equation}
\textit{Significance:} The energy is minimized when $\bmu$ points along $\bB$.  A
misaligned magnetic moment is therefore not in equilibrium: the field tries to rotate it.

\begin{figure}[h]
\centering
\begin{tikzpicture}[>=Stealth, thick, scale=1.0]
  \draw[blue!60!black, thick] (0,0) ellipse (1.4 and 0.55);
  \draw[->, blue!60!black] (-1.35,0.08) -- (-1.0,0.33) node[above left] {$I$};
  \draw[->, red!70!black, line width=1.1pt] (0,0) -- (0.75,1.25)
    node[right] {$\bmu$};
  \foreach \x in {-2.2,-1.4,-0.6,0.2,1.0,1.8,2.6} {
    \draw[->, green!50!black] (\x,-1.3) -- (\x,1.4);
  }
  \node[green!50!black] at (2.9,1.2) {$\bB_{\rm ext}$};
  \draw[dashed] (0,0) -- (0,1.45);
  \draw (0,0.75) arc (90:59:0.75);
  \node at (0.28,0.93) {$\theta$};
\end{tikzpicture}
\caption{A microscopic current loop in a slowly varying external field.  The leading
interaction energy is $U=-\bmu\cdot\bB=-\mu B\cos\theta$.}
\label{fig:lec16-loop-field-energy}
\end{figure}

%------------------------------------------------------------
\subsection{Force on a Magnetic Dipole}\label{subsec:lec16-force}

\subsubsection*{Conceptual Background}
Once the potential energy is known, the force follows from the spatial variation of
that energy.  We keep $\bmu$ fixed while differentiating with respect to the loop's
position; physically, the microscopic current loop is moved through an external field
whose value changes from point to point.

From \eqref{eq:lec16-U-dipole},
\begin{equation}\label{eq:lec16-force-def}
  \bvec{F}=-\nabla U
  =\nabla(\bmu\cdot\bB).
\end{equation}
Thus
\begin{equation}\label{eq:lec16-force}
  \boxed{
    F_i=\partial_i(\mu_j B_j),
    \qquad
    \bvec{F}=\nabla(\bmu\cdot\bB).
  }
\end{equation}
\textit{Significance:} A magnetic dipole feels a translational force only if the
external field is nonuniform.  A uniform field can choose an orientation, but it cannot
choose a preferred position.

For constant $\bmu$ one often writes
\begin{equation}\label{eq:lec16-force-alt}
  \bvec{F}=(\bmu\cdot\nabla)\bB
\end{equation}
when the region is current-free so that $\nabla\times\bB=0$.  The form
\eqref{eq:lec16-force-def} is the direct energy-gradient statement.

\begin{remark}
  This does not contradict the magnetic Lorentz force being perpendicular to the
  instantaneous velocity of a charge.  Holding a steady current loop together and
  moving it as a composite object involves external constraints and induced electric
  fields; the effective dipole energy describes the work needed to move the whole
  object in an inhomogeneous field.
\end{remark}

%------------------------------------------------------------
\subsection{Torque on a Magnetic Dipole}\label{subsec:lec16-torque}

\subsubsection*{Physical Motivation}
In a uniform field, \eqref{eq:lec16-force} gives zero net force.  But the energy
$U=-\mu B\cos\theta$ still depends on orientation, so the dipole must feel a torque.
This is exactly what happens to a compass needle: it does not translate in a uniform
field, but it rotates.

\subsubsection*{Loop Derivation}
For a current element $d\bvec{\ell}$ in a magnetic field,
\begin{equation}\label{eq:lec16-dF-loop}
  d\bvec{F}=\frac{I}{c}\,d\bvec{\ell}\times\bB.
\end{equation}
The torque element about the loop origin is
\begin{equation}\label{eq:lec16-dtau-loop}
  d\boldsymbol{\tau}
  =\br\times d\bvec{F}
  =\frac{I}{c}\,\br\times(d\bvec{\ell}\times\bB).
\end{equation}
Using
\begin{equation}\label{eq:lec16-triple-product}
  \br\times(d\bvec{\ell}\times\bB)
  =d\bvec{\ell}(\br\cdot\bB)-\bB(\br\cdot d\bvec{\ell}),
\end{equation}
and integrating around the loop, one finds that the total derivative terms around
the closed path vanish.  Keeping the leading term for a uniform $\bB$ across the loop,
the result is
\begin{equation}\label{eq:lec16-torque}
  \boxed{\boldsymbol{\tau}=\bmu\times\bB.}
\end{equation}
\textit{Significance:} The torque tries to rotate $\bmu$ toward $\bB$, precisely the
orientation that minimizes $U=-\bmu\cdot\bB$.

The same result follows immediately from the energy:
\begin{equation}\label{eq:lec16-energy-angle}
  U(\theta)=-\mu B\cos\theta,
  \qquad
  \tau_\perp=-\frac{\partial U}{\partial\theta}
  =-\mu B\sin\theta,
\end{equation}
where the sign indicates rotation toward smaller $\theta$.

%------------------------------------------------------------
\subsection{Larmor Precession}\label{subsec:lec16-larmor}

\subsubsection*{Physical Motivation}
If a dipole is tied to an angular momentum, the torque does not simply make it fall
straight into the field direction.  Instead, just like a spinning top in gravity, the
angular momentum vector precesses around the applied field.

For orbital electron motion, Lecture 15 gave
\begin{equation}\label{eq:lec16-mu-L}
  \bmu=\gamma\,\bvec{L},
  \qquad
  \gamma=\frac{e}{2m_ec}.
\end{equation}
The rotational equation of motion is
\begin{equation}\label{eq:lec16-dLdt}
  \frac{d\bvec{L}}{dt}=\boldsymbol{\tau}.
\end{equation}
Using \eqref{eq:lec16-torque},
\begin{equation}\label{eq:lec16-Larmor-equation}
  \frac{d\bvec{L}}{dt}
  =\bmu\times\bB
  =\gamma\,\bvec{L}\times\bB.
\end{equation}
Equivalently,
\begin{equation}\label{eq:lec16-Larmor-box}
  \boxed{
    \frac{d\bvec{L}}{dt}
    =\boldsymbol{\Omega}_L\times\bvec{L},
    \qquad
    \boldsymbol{\Omega}_L=-\gamma\,\bB
    =-\frac{e}{2m_ec}\,\bB .
  }
\end{equation}
\textit{Significance:} The angular momentum keeps its magnitude fixed and rotates
around $\bB$.  The magnetic field creates a preferred axis, not a radial pull.

Indeed,
\begin{equation}\label{eq:lec16-L-magnitude}
  \frac{d}{dt}L^2
  =2\bvec{L}\cdot\frac{d\bvec{L}}{dt}
  =2\gamma\,\bvec{L}\cdot(\bvec{L}\times\bB)=0,
\end{equation}
so $|\bvec{L}|$ is constant.  Also, for uniform constant $\bB$,
\begin{equation}\label{eq:lec16-L-dot-B}
  \frac{d}{dt}(\bvec{L}\cdot\bB)
  =\frac{d\bvec{L}}{dt}\cdot\bB
  =\gamma(\bvec{L}\times\bB)\cdot\bB=0,
\end{equation}
so the angle between $\bvec{L}$ and $\bB$ is constant.  A vector of fixed length and
fixed angle relative to $\bB$ can only precess around $\bB$.

\begin{figure}[h]
\centering
\begin{tikzpicture}[>=Stealth, thick, scale=1.0]
  \draw[->, green!50!black, line width=1.1pt] (0,0) -- (0,3.0)
    node[above] {$\bB$};
  \draw[dashed, gray] (0,1.55) ellipse (1.35 and 0.38);
  \draw[->, red!70!black, line width=1.1pt] (0,0) -- (1.2,1.55)
    node[right] {$\bvec{L},\,\bmu$};
  \draw[->, blue!60!black] (1.35,1.55) arc (0:55:1.35 and 0.38);
  \node[blue!60!black] at (1.65,2.0) {$\boldsymbol{\Omega}_L$};
  \draw (0,0.8) arc (90:52:0.8);
  \node at (0.33,0.95) {$\theta$};
\end{tikzpicture}
\caption{Larmor precession: the torque is perpendicular to $\bvec{L}$, so it changes
the direction of $\bvec{L}$ while preserving its magnitude and its angle with $\bB$.}
\label{fig:lec16-larmor}
\end{figure}

%------------------------------------------------------------
\subsection{From One Dipole to Magnetic Materials}\label{subsec:lec16-materials}

\subsubsection*{Conceptual Background}
The microscopic picture of magnetic matter is a collection of dipoles.  The macroscopic
question is how those dipoles arrange themselves in response to an external field.
This is not fully explainable by classical mechanics: real magnetism, especially spin
and ferromagnetism, requires quantum mechanics.  The classical dipole model still gives
the correct organizing language.

\dfn[dfn:lec16-magnetization]{Magnetization}{
  The magnetization $\bvec{M}$ is magnetic dipole moment per unit volume:
  \begin{equation}\label{eq:lec16-M-def}
    \boxed{\bvec{M}(\br)=\frac{d\bmu_{\rm total}}{dV}.}
  \end{equation}
  It measures the coarse-grained alignment of microscopic magnetic moments.
}

\textit{Significance:} A material produces its own magnetic field once its microscopic
dipoles acquire a nonzero average orientation.

\subsubsection*{Paramagnetism}
In a paramagnet, microscopic moments tend to align with the applied field:
\begin{equation}\label{eq:lec16-paramag}
  \boxed{\bvec{M}\parallel \bB_{\rm ext}.}
\end{equation}
\textit{Significance:} The material reinforces the applied field.  With the field off,
thermal disorder usually randomizes the moments, so the net magnetization disappears.

\subsubsection*{Diamagnetism}
In a diamagnet, the induced magnetic response points opposite the applied field:
\begin{equation}\label{eq:lec16-diamag}
  \boxed{\bvec{M}\ \text{opposes}\ \bB_{\rm ext}.}
\end{equation}
\textit{Significance:} The material partially cancels the applied field.  Classically,
this is connected to the small induced orbital currents created when the external field
perturbs electron motion.

\subsubsection*{Ferromagnetism}
In a ferromagnet, the material can retain magnetization even after the external field
is removed:
\begin{equation}\label{eq:lec16-ferromag}
  \boxed{\bvec{M}\neq 0\quad\text{can remain when}\quad \bB_{\rm ext}=0.}
\end{equation}
\textit{Significance:} A ferromagnet is a material with memory: microscopic moments
remain ordered macroscopically.  Explaining why this order persists is a quantum
many-body problem, not a purely classical magnetostatic calculation.

\subsubsection*{Superconductors as Perfect Diamagnets}
A superconductor responds even more strongly: in the Meissner state the magnetic field
inside is expelled,
\begin{equation}\label{eq:lec16-meissner}
  \boxed{\bB_{\rm inside}=0.}
\end{equation}
\textit{Significance:} A superconductor is the magnetic analogue of an ideal conductor
in electrostatics: it rearranges currents so that the interior field vanishes.  In
magnetism this is the Meissner effect, not merely ordinary perfect conductivity.

\begin{figure}[h]
\centering
\begin{tikzpicture}[>=Stealth, thick, scale=0.95]
  \begin{scope}[xshift=-4.6cm]
    \draw[rounded corners=3pt, fill=blue!5] (-1.2,-0.8) rectangle (1.2,0.8);
    \foreach \x in {-0.75,-0.25,0.25,0.75} {
      \draw[->, red!70!black] (\x,-0.35) -- (\x,0.35);
    }
    \draw[->, green!50!black] (1.55,-0.8) -- (1.55,0.8) node[above] {$\bB_{\rm ext}$};
    \node at (0,-1.25) {\footnotesize paramagnet: aligns};
  \end{scope}
  \begin{scope}
    \draw[rounded corners=3pt, fill=blue!5] (-1.2,-0.8) rectangle (1.2,0.8);
    \foreach \x in {-0.75,-0.25,0.25,0.75} {
      \draw[->, red!70!black] (\x,0.35) -- (\x,-0.35);
    }
    \draw[->, green!50!black] (1.55,-0.8) -- (1.55,0.8) node[above] {$\bB_{\rm ext}$};
    \node at (0,-1.25) {\footnotesize diamagnet: opposes};
  \end{scope}
  \begin{scope}[xshift=4.6cm]
    \draw[rounded corners=3pt, fill=blue!5] (-1.2,-0.8) rectangle (1.2,0.8);
    \foreach \x in {-0.75,-0.25,0.25,0.75} {
      \draw[->, red!70!black] (\x,-0.35) -- (\x,0.35);
    }
    \draw[green!50!black, dashed] (1.55,-0.8) -- (1.55,0.8)
      node[above] {$\bB_{\rm ext}=0$};
    \node at (0,-1.25) {\footnotesize ferromagnet: remembers};
  \end{scope}
\end{tikzpicture}
\caption{Three schematic magnetic responses.  Paramagnets align with the applied
field, diamagnets oppose it, and ferromagnets can retain alignment after the applied
field is removed.}
\label{fig:lec16-material-types}
\end{figure}

%------------------------------------------------------------
\subsection{Summary}\label{subsec:lec16-summary}

\begin{itemize}
  \item The dipole vector potential $\Avec_{\rm dip}=\bmu\times\br/r^3$ gives, for
        $r\neq0$, the dipole field
        $\bB_{\rm dip}=[3\hat{\br}(\bmu\cdot\hat{\br})-\bmu]/r^3$.
  \item For a small current loop in a slowly varying external field,
        $L_{\rm int}=\bmu\cdot\bB$ and $U=-\bmu\cdot\bB$.
  \item The force on a magnetic dipole is $\bvec{F}=\nabla(\bmu\cdot\bB)$; it vanishes
        in a uniform field.
  \item The torque is $\boldsymbol{\tau}=\bmu\times\bB$; a uniform field can rotate
        a dipole even when it cannot translate it.
  \item For orbital electron motion, $\bmu=(e/2m_ec)\bvec{L}$, so
        $d\bvec{L}/dt=(e/2m_ec)\bvec{L}\times\bB$: Larmor precession.
  \item Magnetic materials are understood macroscopically through magnetization
        $\bvec{M}=d\bmu_{\rm total}/dV$.
  \item Paramagnets align with the applied field, diamagnets oppose it, ferromagnets
        retain magnetization, and superconductors expel magnetic field.
\end{itemize}
